\documentclass[multi=tikzpicture,border=9pt]{standalone}
\usepackage[UTF8]{ctex}
\usepackage{amsmath,amssymb}
\usepackage{tikz}
\usetikzlibrary{arrows.meta,calc,patterns,decorations.pathmorphing}
\definecolor{main}{HTML}{17365D}
\definecolor{accent}{HTML}{AC3931}
\tikzset{every picture/.style={font=\small,line width=.75pt},
>={Stealth[length=2mm]},arr/.style={->},box/.style={draw=main,rounded corners=2pt,
fill=main!4,text width=2.7cm,minimum height=.9cm,align=center},
note/.style={align=center,text width=8cm}}
\begin{document}
\begin{tikzpicture}
\fill[main!10] (0,0) circle(1.4);\fill[white] (0,0) circle(.56);\draw (0,0) circle(1.4);\draw (0,0) circle(.56);\draw[dashed,arr] (-5.2,0)--(3.4,0) node[right]{$x$};\fill (0,0) circle(1pt);\node[below] at (0,-.05){$O$};\draw[accent,arr] (-5,0)--(-5,.85) node[above]{物体};\draw[<->] (-5,-2)--(0,-2) node[midway,below]{$u=40\ \mathrm{cm}$};\draw[<->] (0,0)--(1.4,0) node[midway,above]{$R$};\draw[<->] (0,0)--(0,.56) node[midway,right]{$a$};\node[note] at (-.7,-3){同心球面；水 $n=4/3$，空气 $n=1$。中心与水区光路分别处理。};
\end{tikzpicture}
\begin{tikzpicture}
\fill[main!10] (0,0) circle(1.4);\fill[white] (0,0) circle(.56);\draw (0,0) circle(1.4);\draw (0,0) circle(.56);\draw[dashed,arr] (-5.2,0)--(3.4,0) node[right]{$x$};\fill (0,0) circle(1pt);\node[below] at (0,-.05){$O$};\draw[main,arr] (-4.5,.45)--(-1.15,.9)--(1.15,.9)--(3,.45);\draw[accent,arr] (-4.5,.15)--(-1.37,.28)--(-.53,.19)--(.53,.24)--(1.35,.4)--(3,.9);\node[main,above] at (2,1.35){水区光束};\node[accent,below] at (2,-.45){中心光束};\node[note] at (-.5,-2.55){路径为示意，不是精确光线追迹。中心束穿过四界面；水区束绕开气泡。};
\end{tikzpicture}
\begin{tikzpicture}
\draw[arr] (-5.5,0)--(5.7,0) node[right]{$x$};\draw[main] (0,0) circle(1.3);\draw[main] (0,0) circle(.52);\draw[accent,arr] (-5,0)--(-5,1.2);\node[above] at (-5,1.2){物体 $-40$};\draw[accent,dashed,arr] (-1.3,0)--(-1.3,.3);\node[align=center,anchor=north] at (-1.3,-1.55){中心虚像 $-10\ \mathrm{cm}$\\$\mathcal M=+1/4$};\draw[main,arr] (5,0)--(5,-1.2);\node[align=center,anchor=north] at (5,-1.45){水区实像 $+40\ \mathrm{cm}$\\$\mathcal M=-1$};\node[above] at (0,0){$O$};\node[note] at (0,-3){位置相对球心；间距为示意。实虚由最终出射束判断，正倒立由放大率判断。};
\end{tikzpicture}
\begin{tikzpicture}
\node[box] (b0) at (0,0) {外左：$1\to n$\\$\mathcal R=+R$};
\node[box] (b1) at (5,0) {泡左：$n\to1$\\$\mathcal R=+a$};
\node[box] (b2) at (5,-2) {泡右：$1\to n$\\$\mathcal R=-a$};
\node[box] (b3) at (0,-2) {外右：$n\to1$\\$\mathcal R=-R$};
\draw[arr] (b0)--(b1);\draw[arr] (b1)--(b2);\draw[arr] (b2)--(b3);\node[note] at (2.5,-3.35) {三段传播为：水中 $R-a$，空气中 $2a$，水中 $R-a$。用 $(h,n\alpha)$ 统一符号。};
\end{tikzpicture}
\begin{tikzpicture}
\node[box] (b0) at (0,0) {物距先减去 $R$};
\node[box] (b1) at (5,0) {逐界面乘矩阵};
\node[box] (b2) at (5,-2) {$q=-B/D$\\$\mathcal M=A+qC$};
\node[box] (b3) at (0,-2) {球心像坐标 $R+q$};
\draw[arr] (b0)--(b1);\draw[arr] (b1)--(b2);\draw[arr] (b2)--(b3);\node[note] at (2.5,-3.35) {完整传播矩阵与球心等效焦距相互核对；环形水区的球差不在一级模型内。};
\end{tikzpicture}
\end{document}
